\chapter{Background and Literature Review} \label{sec:BackgroundAndLiteratureReview}

\section{Background} \label{sec:Background}

Many stories within video games are experienced in a linear fashion, with the player
having very little autonomy over how the setting progresses. Some games offer you
set choices throughout the playthrough and/or multiple endings, but still ultimately
offer a limited set of experiences. Many players value the replayability of a game,
referring to how much the game can be replayed and exlored before becoming stale,
which is known to be difficult to combine with a strong narrative experience.

By shifting focus away from the story as a whole and onto the characters themselves,
much of the complexity that comes from branching storylines can be avoided. Instead
of considering every character over every narrative path, the characters can be written
to adapt to each path on an individual level, significantly freeing up writer time. This
is referred to as an emergent narrative, as the story ends up inherently emerging from
the actions the player and other characters take. Significant progress has been made
towards being able to represent characters in this manner, but developers still lack a
unified tool that can support this system.

Many tools proposed contain some sort of drama management, which oversees the
overall tension of the story. For example, if events become too dull, the drama
manager might insert some sort of unexpected occurence to make the plot engaging
again. In extreme cases, the drama manager may be in charge of the rising tension,
climax, and resolution of the plot entirely.

\section{Related Work} \label{sec:RelatedWork}

Façade, created in 2003, is often considered one of the earlier attempts at creating an
emergent narrative. Instead of focusing heavily on the characters themselves, Façade
organizes its story into beats, which are pulled by a central drama manager. While
the results are impressive, they are not perfect. Two man years of effort only resulted
in a 20 minute one-act play replayable 6 or 7 times before it was exhausted. The text
input was also notoriously finnicky, which was likely more a result of the time it was
created. \cite{mateas_facade_2003}

Comme il Faut (2011), used to create Prom Week \cite{noauthor_prom_nodate}, saw characters following certain
predefined rules to create the social simulation. \cite{mccoy_comme_2011} This implementation provided
some limitations and was somewhat overly structured, and the project wasn’t publicly
available, leading to the creation of the Ensemble Engine in 2015. \cite{samuel_ensemble_2015} This fixes a
lot of the flexibility issues Comme il Faut had, but still maintains certain structures
that quickly become convoluted. It also hasn’t received support since 2019.

Versu (2012) was another private project that attempted to create believable au-
tonomous characters. It involved a lower-level logic language called Praxis, as well
as a higher-level writer-friendly language named Prompter. \cite{evans_ai_2015} \cite{nelson_prompter_2014} The results from
this project seem promising, but are largely lost to time. The tree and exclusion
language inspired world are powerful, but suffer from a lack of structure that could
make it more susceptible to bugs. Praxish was a recent attempt to recreate the Praxis
language, but it hasn’t been maintained since it was created in 2023 and largely main-
tains the same issues, while not being fully featured. \cite{dameris_praxish_2023} Creating Praxis in a modern
setting and with sturdier foundations will be the primary focus of this paper.

Some other related works include Crosston Tavern (2021) which aims to let charac-
ters take independent action according to goals. \cite{oliver_crosston_2021} Ryan et al. (2015) explore a
system that allows characters to misremember and lie, allowing for even more dra-
matic moments. \cite{ryan_toward_2015} This is an interesting extension, but is outside the scope of
this project. Inform 7, created in 2006, is a popular open-source language for writ-
ing interactive fiction and was a large inspiration for Versu’s higher-level Prompter
language. \cite{inform_7_inform_nodate} ”ink”, created in 2015, is another such system, allowing for narrative
scripting within games using a higher-level writer-friendly language, also supporting
popular game engine integration. \cite{noauthor_ink_nodate} The Sims is a popular video game series that
centers around individual non-player-controlled characters, with the first installment
released in 2000. \cite{arts_sims_2016} Although largely successful, The Sims does not attempt to cre-
ate an overarching narrative with the characters at the center and is more focused on
simulation-style gameplay.
